\documentclass[pdflatex,sn-mathphys-ay]{sn-jnl}

\usepackage{graphicx}
\usepackage{multirow}
\usepackage{amsmath,amssymb,amsfonts}
\usepackage{booktabs}
\usepackage{xcolor}
\usepackage{placeins}

\raggedbottom
\begin{document}

% =============================================================================
\title[Where shade reduces the thermal exposure of active mobility]%
      {Where shade reduces the thermal exposure of active mobility: an
       agent-based simulation study of twelve Spanish cities across five
       climate regimes}
% =============================================================================

\author*[1]{\fnm{Federico} \sur{Bellisardi}}
\email{fbellisardi@ifisc.uib-csic.es}
\author[1]{\fnm{Jos\'e~J.} \sur{Ramasco}}

\affil*[1]{
  \orgdiv{Instituto de F\'isica Interdisciplinar y Sistemas Complejos, IFISC (CSIC-UIB)}, \orgname{CSIC -- Universitat de les Illes Balears}, \orgaddress{\city{Palma}, \country{Spain}}%
}

% =============================================================================
\abstract{%
\textcolor{red}{DRAFT --- results not yet available.} Street-tree canopy and
shade are among the few interventions that directly reduce the thermal
exposure of people walking and cycling, but where to allocate a limited
planting budget is an open, city-specific question, and the benefit of shade
is not unconditional: the same canopy that cools a street in summer can add a
small cold-stress penalty in winter. We study this joint heat-benefit/
cold-cost trade-off in twelve Spanish cities spanning five climate regimes
(coastal Mediterranean, hot semi-arid, continental interior, highland, and
Atlantic oceanic), using an agent-based mobility simulation (NOMAD) driven by
real origin-destination demand (MITMA) on the real road network
(OpenStreetMap), combined with real tree-canopy cover (ESA WorldCover, Urban
Atlas), building height (Urban Atlas, CNIG/IGN PNOA LiDAR), a
reproducible, data-driven calendar of representative days built from the full
2022 hourly thermal-comfort reanalysis (Copernicus UTCI, ERA5-Land), and
in-situ meteorological observations (AEMET) for independent validation. We
compare random, thermal-hotspot, flow-based, active-mobility-weighted,
exposure-optimised and equity-weighted canopy allocation strategies at four
intervention budgets (1\%, 5\%, 10\% and 20\% of street-edge length), evaluate
each against both the cumulative heat dose it avoids and any cold-stress dose
it adds, and report the resulting distribution of net benefit across a range
of relative heat/cold weightings. [Abstract to be completed once results for
all twelve cities are available.]
}

\keywords{urban heat, cold stress, thermal comfort, UTCI, active mobility,
  agent-based simulation, tree canopy, urban greening, equity, climate
  regimes, Spain}

\maketitle

% =============================================================================

Urban heat exposure is not distributed evenly across a city's street network,
and neither is exposure to it: people walking or cycling are exposed to
whatever thermal conditions exist along their route, for the whole duration of
their trip, while people in cars are largely insulated from it. Tree canopy
and shade are among the few interventions that directly reduce pedestrian- and
cyclist-level thermal exposure at the street scale, but where to prioritise
planting is a genuine allocation problem: canopy is costly to establish and
maintain, and a given planting budget can be spent on the streets that are
currently hottest, the streets that carry the most people, the streets that
carry the most active (non-car) travel specifically, or some combination
weighted by who lives nearby and how heat-vulnerable they are. A closely
related remote-sensing study optimised urban greening allocation against
population exposure to land-surface-temperature extremes across 200 cities
worldwide \cite{massaro2023greening}; we address a complementary version of
the same allocation problem using simulated agent-level active-mobility
exposure (not residential population exposure from satellite-derived land
surface temperature) as the objective, which lets us weight streets by who
actually walks or cycles them and when, rather than by who lives nearby.

Two aspects of this allocation problem are, to our knowledge, rarely
addressed together. First, shade is not an unconditional good: the same
canopy that reduces UTCI on a hot afternoon also reduces it, by the same
mechanism, on a cold winter morning, potentially trading a heat benefit for a
cold-stress cost. Evaluating an intervention on a single hot reference day
cannot reveal this trade-off. Second, Spain's own climate diversity --- from
the Atlantic-oceanic north coast to the semi-arid south-east and the
continental interior plateau --- means that a single-city or single-region
study cannot establish how general any given allocation strategy's advantage
is. We address both by (i) building a reproducible, data-driven calendar of
representative days from the full year's hourly UTCI series rather than
picking dates by inspection, explicitly including cold-season and
non-heatwave summer reference days alongside verified heatwave days, and (ii)
applying the same pipeline, unmodified, to twelve Spanish cities chosen to
span five climate regimes. We use real agent-based simulated mobility flows
(not modelled trip generation alone) combined with real thermal-comfort and
canopy data throughout.

% =============================================================================
\section{Results}
\label{sec:results}
% =============================================================================

At the time of writing, the data and simulation pipeline described in
Methods has been validated and run end-to-end for all twelve target cities.
An initial car-only baseline on \texttt{QueueTrafficModel} was run for Palma
de Mallorca (713{,}192 agents; 96.53\% arrival rate, closely matching
NOMAD's own independently published validation figure of 96.5\% for the same
city and demand scale) and, after the threshold recalibration of
Sec.~\ref{sec:stuck}, extended to all twelve cities. The main simulation
campaign then moved to \texttt{LtmTrafficModel} at full demand scale across
the reproducible calendar of representative days (Sec.~\ref{sec:calendar}),
surfacing and fixing a genuine congestion-signal bug in that traffic model's
gate-wait accounting and adding pre-trip rerouting (Sec.~\ref{sec:stuck});
Table~\ref{tab:teleport_ltm} reports the resulting teleport rate for every
city, split by weekday- versus weekend/holiday-representative date, and the
systematic gap between the two is discussed there and in
Sec.~\ref{sec:limitations}.

Per-edge canopy fraction and per-edge, per-hour walk/bicycle flow (with
exact conservation of routed demand, Sec.~\ref{sec:flows}) have been
computed for all twelve cities. The reproducible calendar-construction
procedure (Sec.~\ref{sec:calendar}), applied to the real 2022 hourly UTCI
series for all twelve cities, already yields two findings that inform the
study design rather than its final conclusions: four of five candidate
``ordinary summer'' reference dates supplied for verification turn out to be
locally in the 90th UTCI-max percentile or above (i.e.\ locally extreme) in
at least one city despite falling outside every nationally-declared 2022
heatwave window, and conversely a substantial share (32.5\%) of city-days
inside a nationally-declared heatwave window are \emph{not} locally extreme,
concentrated in the two Atlantic-regime cities (Bilbao, A Coru\~na) --- both
consistent with the motivation in the introduction for verifying reference
dates locally rather than only against a national declaration.

\subsection{Canopy allocation strategy comparison}
\label{sec:results_interventions}

The six canopy-allocation strategies of Sec.~\ref{sec:interventions} have
been compared for ten of the twelve cities at the time of writing (all
except Madrid and Barcelona: Madrid's per-edge canopy/vulnerability/exposure
tables were not yet built at the time of writing despite its car-traffic
simulation being complete, Sec.~\ref{sec:stuck}; Barcelona's canopy and
pedestrian/cyclist flow data were completed most recently and its comparison
was still running), each against a weekday-representative and a
weekend/holiday-representative demand-and-thermal scenario built from the
calendar of Sec.~\ref{sec:calendar}, at four canopy budgets (1\%, 5\%, 10\%,
20\% of street-edge length). Three consistent cross-city patterns emerge.

First, the two flow-aware, exposure-driven strategies (exposure-optimised
and active-mobility-weighted) concentrate addressable heat/cold burden far
more efficiently than the random baseline: at a 5\% length budget, the
random baseline addresses a share of a city's total heat-stress burden
essentially proportional to its length share (as it must, in expectation),
while exposure-optimised and active-mobility-weighted address between 0\%
(in the three heat-negligible cases discussed below) and 86\% of it with the
same 5\% of street length --- a median concentration ratio of roughly 14
relative to random's share among the seventeen combinations where the
comparison is defined (i.e.\ excluding the three heat-negligible cases,
where both random's and the exposure-driven strategies' share are zero) ---
across the 20 city/scenario combinations evaluated (ten cities times two
representative-date groups). Thermal-hotspot allocation
(ranking purely by peak UTCI, with no information about who actually walks
or cycles there) is a considerably weaker proxy for this objective than
either flow-aware strategy, addressing at most 32\% of the same
heat/cold burden at the same budget across every city and scenario tested
--- consistent with this study's motivating premise that where it is hottest
and where people actually experience that heat are not the same question.

Second, equity-weighted allocation trades a documented amount of raw
exposure-efficiency for reaching more socio-demographically vulnerable
populations: across the same 20 city/scenario combinations, its mean
vulnerability index among selected edges exceeds the random baseline's by
0.11--0.51 (median 0.25, on the index's own [0,1] scale), i.e.\ routinely
50--100\% higher vulnerability coverage than an unweighted allocation, at a
real but bounded cost in exposure addressed relative to the pure
exposure-optimised strategy.

Third, and directly testing this study's central heat/cold framing: for the
three city/scenario combinations where ambient UTCI never crosses the
moderate-heat threshold across the whole representative-date group (A
Coru\~na, both scenarios; Bilbao, weekday) --- an Atlantic-oceanic-regime
finding consistent with the climate-regime grouping of
Sec.~\ref{sec:study_areas} --- the heat-only exposure term is exactly zero
everywhere by construction, and exposure-optimised correctly and
automatically falls back to minimising added cold-stress dose instead of
producing an undefined or arbitrary ranking, selecting an allocation with
dramatically lower cold-stress exceedance (by two to five orders of
magnitude) than every other strategy in all three cases. Equity-weighted
allocation, in the same three cases, illustrates a real and informative
tension rather than a second instance of the same fallback: because it
blends the (here, purely cold-driven) net-benefit term with the
vulnerability index at equal weight by construction
(Sec.~\ref{sec:interventions}), it is not guaranteed to also minimise cold
cost, and in one of the three cases (Bilbao, weekday) it selects a
materially \emph{higher}-cold-cost allocation than even the random baseline
because those higher-cold-cost edges happen to serve a more vulnerable
population --- a concrete, quantified instance of the cost side of an
equity-weighted strategy that a heat-only objective could not have surfaced
at all. Together, these confirm that the net heat/cold accounting of
Sec.~\ref{sec:dosebenefit} is doing real work, not degenerating, exactly in
the climate regime where a heat-only objective would have failed silently.

% =============================================================================
\section{Discussion}
\label{sec:discussion}
% =============================================================================

\textcolor{red}{[PENDING.]}

\subsection{Limitations}
\label{sec:limitations}

\textcolor{red}{[DRAFT --- to be expanded as the pipeline matures.]} Shade is
currently approximated as tree-canopy cover only; a genuine building-shadow
model using the now-integrated CNIG/IGN PNOA and Urban Atlas building-height
data (Sec.~\ref{sec:building_height}) is planned but not yet implemented, and
in the interim a simplified evergreen/deciduous seasonal-canopy comparison
(Sec.~\ref{sec:phenology}) uses illustrative, not species- or city-validated,
leaf-on/leaf-off timing. The magnitude of canopy-shade UTCI reduction is a
documented literature-range assumption, not a value calibrated against
in-situ shaded/unshaded UTCI measurements for these specific cities. The
heat/cold net-benefit trade-off (Sec.~\ref{sec:dosebenefit}) depends on a
relative-weight parameter $\lambda$ that this study does not have grounds to
fix on its own behalf; results are reported across a range of $\lambda$
values rather than at a single chosen point. Barcelona's road network was,
for most of this study, blocked by a genuine infinite-loop defect in
NOMAD's own network-simplification preprocessing step (repeated
re-contraction of already-absorbed nodes, triggered by Barcelona's specific
topology); this has since been diagnosed and fixed with an authorised,
tested change to that preprocessing code (not to any simulation logic), and
Barcelona's network now builds in seconds rather than hanging, but its
canopy, pedestrian/cyclist flow and intervention-strategy comparison were
completed most recently among the twelve cities and are reported in
Sec.~\ref{sec:results_interventions} only to the extent available at the
time of writing. Madrid's origin-destination matrix is roughly twenty times
larger than Palma de Mallorca's baseline, with a proportionally larger
simulation cost; its car-traffic simulation is complete for the full
calendar (Table~\ref{tab:teleport_ltm}), but its per-edge canopy,
vulnerability and exposure tables needed for the intervention-strategy
comparison were not yet built at the time of writing. NOMAD has no implemented
public-transit model, so transit-mode MITMA trips are excluded from
simulation entirely rather than reassigned to another mode. NOMAD's native
output instrumentation is car-only; pedestrian and cyclist flow relies on an
offline-routing reconstruction (Sec.~\ref{sec:flows}) that is exact under
NOMAD's own no-congestion-feedback design for those modes, but has not been
cross-validated against a full discrete-event run with equivalent per-mode
instrumentation, since NOMAD does not currently expose one. The
socio-demographic vulnerability index currently has full census-section
coverage only within each city's core municipality, not its whole Functional
Urban Area; edges outside that coverage fall back to the FUA-wide mean
vulnerability rather than a locally estimated one. Eight of the twelve
cities currently use the coarser, MITMA-zone-resolution, age-only
vulnerability proxy rather than the INE income/poverty variant
(Sec.~\ref{sec:vulnerability}), pending manual portal export of the
remaining cities' income/poverty tables. The ``flow'' strategy's all-mode
ranking (Sec.~\ref{sec:interventions}) is genuinely all-mode only where a
car-traffic simulation has already been exported into the same per-edge
table as pedestrian/cyclist flow; for cities where that join has not yet
been performed, ``flow'' and ``active-mobility'' currently rank identically
(both walk+cycle flow only), a documented, not silently absorbed,
simplification for those cities. The systematic weekday-versus-weekend/
holiday teleport-rate gap of Table~\ref{tab:teleport_ltm} --- present in
every city at a broadly similar order of magnitude despite very different
city sizes and network topologies, which argues against a
city-specific network or routing defect and toward a property of the
weekend/holiday origin-destination matrix construction itself --- is
reported here as an open question rather than a resolved one; identifying
its specific mechanism is ongoing work.

% =============================================================================
\section{Conclusion}
\label{sec:conclusion}
% =============================================================================

\textcolor{red}{[PENDING --- to be written once results are available for all
twelve cities.]}

% =============================================================================
\section{Methods}
\label{sec:methods}
% =============================================================================

Table~\ref{tab:data_sources} catalogues every real, third-party data source
used in this study, its provider, its role in the pipeline, and its current
integration status; it is kept as the single authoritative summary and
updated whenever a new source is added. Each source is described in full,
with its specific role and any documented simplifying assumptions, in the
corresponding subsection below.

\begin{table}[h]
\caption{Data sources used in this study.}
\label{tab:data_sources}
\begin{tabular}{@{}p{2.6cm}p{3.3cm}p{3.6cm}p{1.6cm}@{}}
\toprule
Source & Provider & Role & Status \\
\midrule
OpenStreetMap & OSM contributors / Geofabrik & Road network topology & Integrated (11/12) \\
GHSL Functional Urban Areas & JRC / OECD & Study-area boundaries & Integrated (12/12) \\
MITMA mobility study & Ministerio de Transportes (Spain) & Origin-destination travel demand & Integrated (11/12) \\
ESA WorldCover 2021 v200 & European Space Agency & Tree canopy cover & Integrated (11/12) \\
ERA5-Land & Copernicus C3S & Meteorological reanalysis & Integrated \\
Derived UTCI time series (full year) & Copernicus CDS & Ambient thermal comfort, calendar construction & Integrated (12/12) \\
AEMET OpenData & AEMET (Spain) & In-situ met.\ cross-check & Integrated (12/12) \\
Atlas de Renta de los Hogares & INE (Spain) & Vulnerability weighting & Integrated (municipality-level) \\
Urban Atlas (LC/LU, Street Tree Layer, Building Height) & Copernicus Land Monitoring Service & Land use, canopy cross-check, building height & Integrated (12/12) \\
PNOA LiDAR nDSM & CNIG / IGN (Spain) & Building/vegetation shadow & Integrated (12/12) \\
\botrule
\end{tabular}
\end{table}

\subsection{Study areas}
\label{sec:study_areas}

Twelve cities --- Palma de Mallorca, Zaragoza, Murcia, Valladolid, Madrid,
Barcelona, Valencia, Sevilla, C\'ordoba, Granada, Bilbao and A Coru\~na ---
were selected to span both a range of city sizes and five Spanish climate
regimes: coastal Mediterranean (Barcelona, Valencia, Palma de Mallorca),
hot/semi-arid (Murcia, Sevilla), continental interior (Zaragoza, Valladolid,
Madrid, C\'ordoba), highland (Granada), and Atlantic oceanic (Bilbao, A
Coru\~na). This climate-regime grouping is a provisional working
classification based on general geographic knowledge and is intended to be
confirmed against a citable K\"oppen--Geiger dataset before being relied upon
for any quantitative claim. Each city's study area is its FUA, not its
municipal boundary, following the Global Human Settlement Layer (GHSL)
Functional Urban Areas delineation \cite{ghsl_fua}, sourced from the JRC/OECD
GHSL data package (layer \texttt{FUAs}, matched by FUA name; see
\texttt{network\_prep/boundaries.py} in the code repository). All twelve
cities now have a validated routing network, mobility demand and completed
car-traffic simulation (Table~\ref{tab:teleport_ltm}); the per-edge canopy,
vulnerability and pedestrian/cyclist-flow tables needed for the
intervention-strategy comparison remain incomplete for Madrid and Barcelona
at the time of writing (Sec.~\ref{sec:limitations}).

\subsection{Simulation engine}
\label{sec:nomad}

Agent-based mobility simulation is performed by NOMAD (Network-based Open
Mobility Agent Dynamics) \cite{nomad2026}, a discrete-event agent-based
mobility simulator implemented in C++20, used here strictly as a
black-box engine: its simulation logic (routing, traffic dynamics, event
scheduling) is treated as fixed and is not modified for this study, except for
one narrowly-scoped, tested addition to its Python bindings (exposing an
existing C++ router/traffic-model attachment API and a deterministic
per-edge, per-mode travel-time query, neither of which alters any simulation
behaviour) needed to extract per-mode flow that NOMAD's own instrumentation
does not otherwise expose (Sec.~\ref{sec:flows}). NOMAD supports car,
pedestrian and bicycle agents on a shared road-network graph derived from
OpenStreetMap, with mode-aware routing (via A*) and road-class accessibility
per mode, and a BPR volume-delay queueing traffic model for car congestion.

\subsection{Stuck-agent handling and its calibration}
\label{sec:stuck}

An agent whose elapsed trip time exceeds a configurable multiple of its
route's free-flow time (\texttt{stuck\_threshold\_ratio}), or an absolute
hard cap (\texttt{stuck\_max\_hours}), whichever binds first, is removed from
the simulation (``teleported'') rather than left to run indefinitely; this
mechanism is identical across NOMAD's two available traffic models
(\texttt{QueueTrafficModel}, a BPR volume-delay queue, and
\texttt{LtmTrafficModel}, a Daganzo 1994 single-cell-per-link Cell
Transmission Model with genuine receiving-side spillback), since both are
invoked through the same \texttt{ITrafficModel} interface and neither
implements this bookkeeping itself. An initial single-day, car-only pilot on
\texttt{QueueTrafficModel} surfaced a real calibration issue in this
project's own wrapper (not in NOMAD's simulation logic): the values this
project's code was passing (\texttt{stuck\_threshold\_ratio=20},
\texttt{stuck\_max\_hours=4}) had, in effect, only ever been exercised
against Palma de Mallorca before this study's twelve-city extension, and did
not generalise --- an otherwise identical car-only baseline for Zaragoza
showed a 23.14\% teleport rate at these values, against 1.28\% for Palma.
Relaxing the thresholds to a more moderate (50, 8h) pair, adopted as this
project's new default, reduced this to near zero for every city tested under
that same single-day, off-peak-scale pilot.

That pilot, however, used \texttt{QueueTrafficModel} and a single reference
day at reduced demand scale, and does not represent this study's final
results. The main simulation campaign uses \texttt{LtmTrafficModel}
throughout, run at full demand scale across the reproducible calendar of
representative days (Sec.~\ref{sec:calendar}), because it is the traffic
model capable of genuine spillback (a downstream link at capacity blocks
upstream discharge, rather than a car simply queueing invisibly past a full
link as under the BPR volume-delay formulation) --- the congestion signal the
canopy-allocation objective of Sec.~\ref{sec:dosebenefit} needs to be
realistic. Scaling \texttt{LtmTrafficModel} to full demand surfaced a second,
distinct issue, this time in the traffic model's congestion-signal
bookkeeping rather than in this project's threshold calibration: the model's
BPR-style travel-time term saturates once a link's occupancy reaches its
storage capacity (\texttt{on\_enter} refuses further entries once full), so
the reported travel-time ratio plateaus at a fixed ceiling ($\times 1.15$
in this implementation) no matter how many additional hours a fully-gridlocked
link remains blocked --- confirmed directly on a real Palma de Mallorca
corridor edge that sat at $\sim$100\% storage occupancy for over sixteen
continuous hours while its reported travel-time ratio never moved. A
gate-wait counter already existed in the codebase to compensate for exactly
this blind spot, but was wired only to an opt-in discharge-capacity
mechanism, not to the spillback path that gates every link's exit
unconditionally; extending it to track continuous blocking from either cause
(spillback or discharge-capacity exhaustion) fixed the signal without
touching the underlying Cell Transmission Model dynamics themselves, and
reduced Palma de Mallorca's full-day teleport rate from 39.5\% to 24.1\% in
isolation. A complementary fix addresses a separate failure mode: every
agent's route is computed once, at simulation start, against free-flow
costs, so an agent still waiting to depart hours later has no way to react
to congestion that has since formed along its route; a bulk, parallel
pre-trip rerouting pass, applied to each waiting agent shortly before its own
scheduled departure using the same congestion-aware router already used for
in-trip rerouting, closes this gap without the prohibitive cost of routing
synchronously at every individual departure. With both fixes applied, the
full twelve-city, five-date-calendar campaign (Table~\ref{tab:teleport_ltm})
shows a persisting, systematic weekday/weekend-or-holiday teleport-rate gap
--- median 1.16\% on weekday dates versus 20.47\% on the calendar's two
verified summer/holiday dates, roughly an order of magnitude higher, in
every city without exception --- which we report here as a genuine finding
about demand-driven congestion under this study's OD matrices rather than as
a residual simulation artifact still to be eliminated; its root cause (most
plausibly a property of the weekend/holiday OD matrix specifically, since the
gap recurs at a similar order of magnitude across cities of very different
size and topology) is not yet established and is noted as an open question
in Sec.~\ref{sec:limitations}.

\begin{table}[h]
\caption{Agent teleport rate under the final \texttt{LtmTrafficModel}
configuration (spillback gate-wait fix, Sec.~\ref{sec:stuck}, plus
pre-trip rerouting), full demand scale, weekday-representative vs.\
weekend/holiday-representative dates of the reproducible calendar
(Sec.~\ref{sec:calendar}).}
\label{tab:teleport_ltm}
\begin{tabular}{@{}lcc@{}}
\toprule
City & Teleport rate (weekday) & Teleport rate (weekend/holiday) \\
\midrule
Palma de Mallorca & 4.80\% & 25.21\% \\
Zaragoza          & 0.42\% & 13.61\% \\
Murcia            & 0.11\% & 3.45\% \\
Valladolid        & 0.00\% & 0.02\% \\
Madrid            & 6.37\% & 36.56\% \\
Valencia          & 0.86\% & 23.87\% \\
Sevilla           & 6.65\% & 33.02\% \\
C\'ordoba         & 0.63\% & 10.00\% \\
Granada           & 0.28\% & 8.41\% \\
Bilbao            & 2.04\% & 20.47\% \\
A Coru\~na        & 4.24\% & 25.04\% \\
Barcelona         & 1.16\% & 17.16\% \\
\midrule
Median (all 12)   & \multicolumn{2}{c}{1.16\% (weekday) / 20.47\% (weekend/holiday)} \\
\botrule
\end{tabular}
\end{table}

\subsection{Road network}
\label{sec:network}

Road network topology for each city is extracted from OpenStreetMap
\cite{openstreetmap}, via regional \texttt{.osm.pbf} extracts published by
Geofabrik GmbH, clipped to each city's FUA polygon and simplified (degree-2
node contraction) by NOMAD's own \texttt{simplify\_osm.py} preprocessing
script, called as a subprocess and not reimplemented (see
\texttt{nomad\_wrapper/build.py} in the code repository), with one
authorised bug fix to that preprocessing step's \texttt{NetworkCleaner}
component (Sec.~\ref{sec:limitations}): it re-contracted already-absorbed
nodes indefinitely on Barcelona's specific topology, an infinite loop rather
than a slow computation, fixed by excluding nodes already absorbed by a
prior contraction pass from being re-selected. For Palma de Mallorca this
yields a routing graph of 79{,}111 nodes and 186{,}414 directed edges after
simplification; graph sizes for the other eleven cities range from
26{,}471 nodes / 67{,}120 edges (C\'ordoba) to 495{,}697 nodes /
1{,}268{,}636 edges (Madrid); Barcelona's now-buildable network is the
second largest of the twelve at 917{,}528 edges.

\subsection{Mobility demand}
\label{sec:demand}

Origin-destination travel demand is derived from Spain's Ministry of
Transport (Ministerio de Transportes y Movilidad Sostenible, MITMA)
mobile-phone-derived mobility study, publicly released as anonymised,
aggregated daily zone-to-zone trip matrices
\cite{mitma_movilidad}. Zone-level trips are converted to node-level,
per-mode origin-destination pairs by NOMAD's own \texttt{build\_od.py}
script (unmodified), which additionally restricts each mode's candidate
origin/destination nodes to that mode's own largest strongly-connected
component of the road network, so that every generated trip is guaranteed
routable for its assigned mode, and which generates both an average-weekday
and an average-weekend/holiday matrix (weekend/holiday classification via the
same national-holiday calendar for Spain used throughout). Car mode trip
counts are converted from MITMA person-trips to vehicle-trips using an
average vehicle-occupancy factor of 1.20 persons per car (Spain average). For
Palma de Mallorca, February 2022 average-weekday demand totals 1{,}922{,}543
person-trips across car (713{,}192), walk (573{,}925), bicycle (82{,}594) and
public transit (552{,}832, excluded from simulation --- NOMAD has no
implemented transit model at time of writing).

\subsection{Per-mode flow extraction}
\label{sec:flows}

NOMAD's own output instrumentation (its GeoJSON snapshot writer and
in-memory link-state accessor) records occupancy, inflow and travel time only
for car-mode agents: this is a property of the underlying traffic model, which
is invoked only for car agents (pedestrians and cyclists never contend for
car-calibrated link capacity, and NOMAD implements no pedestrian/bicycle
congestion model), confirmed by direct inspection of the simulation engine's
source. Car-mode per-edge, 15-minute-bin flow, occupancy, mean travel time,
free-flow travel time and derived congestion metrics are therefore extracted
directly from NOMAD's own periodic network-state snapshots. Pedestrian and
bicycle per-edge, per-hour flow, person-seconds and travel time are instead
obtained by direct routing (NOMAD's A* router, mode-aware, queried through
its Python bindings) rather than by running the discrete-event simulation:
because neither mode ever experiences congestion feedback in NOMAD, each
agent's pre-computed route is deterministically its entire realised trip, so
its per-edge timing can be computed exactly from the route and each mode's
free-flow speed, without the computational cost of discrete-event simulation.
Per-agent departure times, sampled uniformly within each origin-destination
row's MITMA-derived hourly window, are propagated to per-edge, per-bin
expected flow and dwell-time (person-seconds) contributions in closed form
(exact overlap integrals of the departure window against each time bin, not
Monte Carlo sampling), so that the resulting per-edge-per-bin tables are
exactly reproducible. Conservation of total routed demand (per-edge flow at
trip origin equal to total routed agents; total person-seconds equal to the
sum of in-simulation-window and next-day-spillover contributions) is checked
programmatically for every extraction (see
\texttt{nomad\_wrapper/walk\_bike\_routes.py}), and has been confirmed exact
for every city processed to date.

\subsection{Tree canopy}
\label{sec:canopy}

Tree canopy cover is obtained from the European Space Agency's WorldCover
2021 v200 global land-cover product \cite{esa_worldcover}, a 10\,m-resolution
global land-cover classification derived from Sentinel-1 and Sentinel-2
satellite imagery. Canopy fraction per road-network edge is computed as the
fraction of ``tree cover'' (WorldCover class~10) pixels within a 15\,m buffer
of each edge's centreline (straight line between its endpoint nodes, since the
simplified routing graph does not retain intermediate street geometry), using
zonal statistics against the relevant WorldCover tile(s) reprojected to each
city's local UTM zone (29N, 30N or 31N depending on longitude) for metric
buffering. Mean edge-level canopy fraction varies substantially across the
eleven cities processed to date, from 14.4--14.6\% in the two semi-arid/
continental cities with the least canopy (Zaragoza, Murcia) to 25.5\% in the
Atlantic-regime city with the most (Bilbao), consistent with the climate-
regime grouping of Sec.~\ref{sec:study_areas}. The Copernicus Urban Atlas
Street Tree Layer (Sec.~\ref{sec:building_height}) provides a second,
independent per-city canopy source, now integrated but not yet
cross-validated against the WorldCover-derived fraction used above.

\subsection{Meteorological reanalysis and thermal comfort}
\label{sec:thermal}

Ambient (unshaded-baseline) thermal comfort is quantified using the
Universal Thermal Climate Index (UTCI) \cite{brode2012utci}, obtained
directly as a derived reanalysis product from the Copernicus Climate Data
Store's \texttt{derived-utci-historical-timeseries} dataset
\cite{cds_utci}, rather than recomputed from first principles, since the CDS
product already applies the official UTCI methodology to ERA5 reanalysis
inputs. The complete hourly UTCI series for calendar year 2022 (8{,}760
hourly values) has been retrieved for every one of the twelve cities in a
single request each, rather than one request per reference day, after
confirming directly against the live API that it accepts a genuine date-range
request. Complementary raw meteorological fields (2\,m temperature, 2\,m
dewpoint temperature, 10\,m wind components, surface pressure, downward solar
and thermal radiation) are obtained from the Copernicus \texttt{ERA5-Land}
hourly reanalysis \cite{era5_land} for cross-validation and for driving a
from-first-principles UTCI computation (via the \texttt{pythermalcomfort}
implementation of the Bröde et al.\ \cite{brode2012utci} polynomial
regression) where a real UTCI product is not available for a given
city/date/scenario combination.

Street-level shade is approximated, for this version of the study, as tree
canopy cover only (Sec.~\ref{sec:canopy}); building-shadow casting from the
now-integrated building-height data (Sec.~\ref{sec:building_height}) is
planned but not yet implemented (Sec.~\ref{sec:limitations}). The ambient
UTCI is adjusted per road edge and hour by a canopy-shade correction, gated by
actual solar elevation (computed from a standard closed-form solar-position
equation for each city's latitude/longitude and simulation date, not from
ERA5's accumulated solar-radiation field, whose accumulation-reset convention
at the 00~UTC time step introduces an ambiguity we chose not to rely on) so
that the correction is zero at night and scales with daylight strength through
the day. The magnitude of full-canopy-shade UTCI reduction at solar noon
(default $6\,^\circ$C) is an explicit, documented, calibratable assumption
drawn from the range reported in the urban-microclimate literature for UTCI
specifically (as opposed to mean-radiant-temperature reductions alone, which
are typically reported as larger), not a value measured or fitted for these
twelve cities.

\subsection{Reproducible thermal calendar}
\label{sec:calendar}

Rather than selecting reference dates by inspection, we build a reproducible
calendar of representative days directly from each city's full 2022 hourly
UTCI series. For each calendar month, we identify the day whose daily maximum
UTCI is closest to that month's median, either independently per city
(``per-city medoid'') or, as our primary, recommended choice, as a single
common date per month that minimises
$\sum_c |\mathrm{UTCI}_{\max}(c,d) - \mathrm{median}_m(\mathrm{UTCI}_{\max,c})| / \mathrm{IQR}_m(c)$
across all twelve cities $c$ for candidate day $d$ in month $m$ --- directly
comparable across cities since every city is evaluated on the same calendar
date, at the cost of being individually less representative for any single
city than that city's own medoid; the per-city variant is retained as a
robustness check. This yields twelve monthly reference dates for 2022. We
additionally verify fifteen candidate dates supplied for the specific
purpose of contrasting ordinary and extreme summer conditions --- five
``ordinary summer'' dates and ten dates drawn from the two heatwaves
officially declared by AEMET for summer 2022 (9--26 July and 30 July--15
August) --- not by trusting the national declaration alone, but by computing
each city-day's UTCI-max percentile within that city's own 2022 distribution,
its meteorological season, and, for the minimum daily UTCI, an independent
cold-stress category using the same official UTCI stress-category thresholds
used throughout this study. The union of the twelve monthly dates and the
fifteen verified dates, after removing one exact duplicate, gives a final
26-date calendar (Sec.~\ref{sec:results}); the minimum date subset needed for
a main-text-scale analysis and the full 26-date set for a robustness
appendix are both reported from the same underlying classification table.

\subsection{In-situ meteorological observations}
\label{sec:aemet}

Daily station-level meteorological observations are obtained from AEMET
(Agencia Estatal de Meteorolog\'ia) OpenData \cite{aemet_opendata}, Spain's
national meteorological agency, for one representative station per city,
selected as the station nearest each city's true urban-centre landmark point
(not the geometric centroid of its FUA polygon, which for irregularly-shaped
FUAs can be materially offset from the city itself --- confirmed for two
cities in this study, where the nearest-by-centroid station would have been
several kilometres further from the actual urban core than the
nearest-by-true-centre station). These observations are used as an
independent, real-world cross-check against the ERA5-Land/UTCI reanalysis
products described above, not as a substitute input to the simulation
pipeline.

\subsection{Socio-demographic vulnerability}
\label{sec:vulnerability}

Heat-vulnerability weighting for the equity-oriented intervention strategy
(Sec.~\ref{sec:interventions}) uses Spain's National Statistics Institute
(Instituto Nacional de Estad\'istica, INE) Atlas de Distribuci\'on de Renta de
los Hogares \cite{ine_adrh}, a census-section-level household income and
poverty-threshold indicator dataset, spatially joined to census-section
boundaries retrieved live from INE's own OGC API Features geometry service.
Two real issues were found and corrected while integrating this source: the
service's response mixes true section-level polygons with one
district-aggregate polygon per district that geometrically overlaps all of
that district's real sections, which is filtered out before the spatial
join; and a flat single request silently truncates municipalities with more
than 2{,}000 census sections (confirmed for Madrid, whose 2{,}483 real
sections were being truncated to 2{,}000), fixed with standard OGC API
Features pagination. At the time of writing, census-section coverage is
available only for each city's core municipality rather than its whole FUA
(Sec.~\ref{sec:limitations}).

The Atlas de Renta's income and poverty-threshold indicator tables
themselves, unlike its census-section boundaries, are retrieved from INE's
table-query portal as a manual per-municipality export rather than a
queryable API, and were obtained this way for only four of the twelve
cities (Palma de Mallorca, Zaragoza, Murcia, Valladolid) at the time of
writing. For the remaining eight, we substitute a coarser but fully
automated proxy: each city's population aged 65 and over as a share of
total resident population, per MITMA ``distrito'' zone (the same official
zonification and population reference already used for origin-destination
demand, Sec.~\ref{sec:demand}, joined to one day's age-stratified
``personas'' file from the same source), spatially assigned to road-network
edges by the same zone-to-edge join method as the income/poverty variant.
Age-65-plus share is itself one of the three indicators the underlying
composite-index formula is designed to accept (Sec.~\ref{sec:interventions}
uses it alone, at full weight, where income/poverty data is unavailable);
MITMA's zones are coarser than INE's census sections, so this variant
resolves vulnerability at zone rather than section granularity. This is a
real, documented resolution and indicator-composition difference between
the two groups of cities, not a value judgement about which population is
more vulnerable, and is reported per city alongside the intervention
results in Sec.~\ref{sec:results_interventions}.

\subsection{Building height and detailed land cover}
\label{sec:building_height}

Two additional real data sources complement the canopy-only shade
approximation of Sec.~\ref{sec:canopy}: the Copernicus Land Monitoring
Service's Urban Atlas \cite{urban_atlas}, a high-resolution
($\leq$0.25~ha minimum mapping unit) land-use/land-cover vector product for
European Functional Urban Areas, together with its per-FUA Street Tree Layer
and Building Block Height raster products, retrieved via the Copernicus Data
Space Ecosystem's S3 object store using a live per-product catalogue (rather
than hardcoding file names, since each product's version/date suffix is
reprocessed over time); and Spain's national mapping agency (Centro Nacional
de Informaci\'on Geogr\'afica, CNIG/IGN) second-coverage PNOA airborne
LiDAR-derived normalised digital surface models for buildings and vegetation
\cite{cnig_pnoa_lidar}, at 2.5\,m grid spacing, for which the exact set of
national 1:50{,}000 map sheets intersecting each city's FUA was computed from
the official IGN sheet grid rather than assumed from a fixed neighbourhood
around each city's centre --- a real gap this correction closed: the fixed-
neighbourhood approach previously in use under-covered even the four
original study cities, most severely Palma de Mallorca, whose FUA extends
across most of Mallorca island. A previously-documented limitation, that no
usable building-height product could be found for these cities, is now known
to have been a mis-scoped single national file (covering only the Canary
Islands) rather than a genuine data-availability gap: the correct, per-FUA
Urban Atlas Building Height product is available for all twelve cities.
Both height products would allow a genuine building-shadow-casting shade
model rather than the canopy-only approximation currently in use; this
remains planned, not yet implemented (Sec.~\ref{sec:limitations}).

\subsection{Thermal dose and shade benefit/cost accounting}
\label{sec:dosebenefit}

For a given canopy scenario, city, and representative day, we define the
cumulative heat and cold thermal dose experienced by active-mode travellers as
\[
  \mathrm{heat\_dose} = \sum_{e,t} N[e,t]\,\max(\mathrm{UTCI}[e,t]-\theta_{\mathrm{heat}},0)\,\Delta t,
  \qquad
  \mathrm{cold\_dose} = \sum_{e,t} N[e,t]\,\max(\theta_{\mathrm{cold}}-\mathrm{UTCI}[e,t],0)\,\Delta t,
\]
where $N[e,t]$ is the walking-plus-cycling person-count on edge $e$ during
hour $t$ (Sec.~\ref{sec:flows}) and $\theta_{\mathrm{heat}}=26\,^\circ$C,
$\theta_{\mathrm{cold}}=9\,^\circ$C reuse the same official UTCI
stress-category boundaries used throughout this study (onset of ``moderate
heat'' and ``slight cold'' stress respectively), rather than newly chosen
thresholds. Shade's heat benefit and cold cost are then
\[
  \mathrm{heat\_benefit} = \mathrm{heat\_dose}_{\mathrm{ambient}} - \mathrm{heat\_dose}_{\mathrm{shaded}},
  \qquad
  \mathrm{cold\_cost} = \mathrm{cold\_dose}_{\mathrm{shaded}} - \mathrm{cold\_dose}_{\mathrm{ambient}},
\]
and net benefit for a given relative weight $\lambda$ is
$\mathrm{net\_benefit}(\lambda) = \mathrm{heat\_benefit} - \lambda\,\mathrm{cold\_cost}$.
Because $\lambda$ encodes a value judgement this study has no independent
grounds to fix, results are reported across a range of $\lambda$ values
(Sec.~\ref{sec:limitations}) and the two un-combined objectives
(heat\_benefit, cold\_cost) are always reported alongside any $\lambda$-combined
figure. This accounting has been validated end-to-end against real data for
Palma de Mallorca: on 2022-02-08 (winter), an evergreen-canopy scenario shows
zero heat benefit (as expected, ambient UTCI never approaches
$\theta_{\mathrm{heat}}$ in February) and a small positive cold cost, while a
deciduous-canopy scenario (Sec.~\ref{sec:phenology}) shows zero cold cost on
the same date, since a dormant deciduous canopy casts no shade to remove; on
2022-07-14 and 2022-07-20 (both inside the officially-declared July
heatwave), both canopy types show an identical, substantial heat benefit and
zero cold cost, as expected since a deciduous canopy is modelled as fully
leafed in mid-summer.

\subsection{Canopy phenology: evergreen versus deciduous}
\label{sec:phenology}

To assess whether the heat-benefit/cold-cost trade-off depends on planting
evergreen versus deciduous species, we implement three canopy-type scenarios:
a year-round (``evergreen'') shade fraction, a seasonally-attenuated
(``deciduous'') shade fraction that ramps linearly to zero over a
configurable window around leaf-off and back to full shade around leaf-on,
and an explicit ``no distinction'' scenario numerically identical to
evergreen but naming the assumption. Real phenology literature exists for
Mediterranean urban trees (e.g.\ studies of urbanisation effects on tree
phenology in Madrid, and Platanus flowering-model studies for other Spanish
cities), but extracting a single verified leaf-on/leaf-off day-of-year per
city and species from that literature would require reading the full papers
in depth rather than a title or abstract, which we have not yet done;
consequently the current leaf-on/leaf-off timing (day-of-year 100 and 305,
respectively, with a 20-day linear ramp) is an explicit, round-number
placeholder for sensitivity analysis, not a value we claim is measured or
validated for any of these twelve cities or their actually-planted species.

\subsection{Intervention strategies}
\label{sec:interventions}

Six canopy-allocation ranking strategies are compared --- random (null
baseline), thermal-hotspot (highest peak UTCI first), flow (highest total
all-mode flow first, where car-mode flow data is available;
Sec.~\ref{sec:results_interventions}), active-mobility (highest
pedestrian+cyclist flow first), exposure-optimised and equity-weighted --- at
four canopy budgets defined as a share of total street-edge length: 1\%,
5\%, 10\% and 20\%. Exposure-optimised ranks edges by net heat/cold benefit,
$\mathrm{exposure\_score} - \lambda\,\mathrm{cold\_exposure\_score}$, the
per-edge analogue of Sec.~\ref{sec:dosebenefit}'s heat\_dose/cold\_dose
applied to each edge's actual (already-shaded, not ambient-vs-shaded)
current thermal exposure rather than to shade's marginal benefit --- a
deliberate choice, since an ambient-vs-shaded attribution is exactly zero by
construction on any edge with little or no canopy today, precisely the
edges a ``where to plant next'' ranking most needs to surface. At
$\lambda=1$ (the symmetric, equal-weight point already present in the
$\lambda$-sweep of Sec.~\ref{sec:dosebenefit} and used here as the ranking
default, not as a claim that this is the correct value), this reduces to
exposure\_score alone whenever cold\_exposure\_score is negligible, and
automatically pivots toward minimising cold\_exposure\_score in the reverse
case (Sec.~\ref{sec:results_interventions}). Equity-weighted allocation
blends min-max-normalised net benefit and the vulnerability index of
Sec.~\ref{sec:vulnerability} at equal (0.5/0.5) weight. Each strategy/budget
combination is evaluated once per city against a weekday-representative and
a weekend/holiday-representative scenario, each an average over the
respective dates of the reproducible calendar (Sec.~\ref{sec:calendar}) ---
a coarser two-scenario structure than the full per-date calendar, adopted
for the results reported here for computational tractability across twelve
cities; the full per-date and per-canopy-phenology-scenario
(Sec.~\ref{sec:phenology}) evaluation remains future work
(Sec.~\ref{sec:limitations}).

\backmatter

\section*{Abbreviations}
AEMET: Agencia Estatal de Meteorolog\'ia (Spanish State Meteorological
Agency); CDS: Climate Data Store; CDSE: Copernicus Data Space Ecosystem; FUA:
Functional Urban Area; GHSL: Global Human Settlement Layer; IGN: Instituto
Geogr\'afico Nacional; INE: Instituto Nacional de Estad\'istica; MITMA:
Ministerio de Transportes y Movilidad Sostenible; OSM: OpenStreetMap; PNOA:
Plan Nacional de Ortofotograf\'ia A\'erea; UTCI: Universal Thermal Climate
Index.

\bmhead{Acknowledgements}
This work uses NOMAD, an agent-based mobility simulator developed by
F.B.\ \cite{nomad2026}. Mobility demand data were obtained from MITMA's public
mobility study; road network data from OpenStreetMap contributors under the
Open Database Licence; land-cover and building-height data from ESA
WorldCover and Copernicus Urban Atlas; meteorological reanalysis and
thermal-comfort data from the Copernicus Climate Data Store; in-situ
meteorological data from AEMET OpenData; socio-demographic indicators from
INE; building/vegetation height from CNIG/IGN PNOA LiDAR.

\section*{Declarations}

\bmhead{Funding}
% AUTHOR INPUT REQUIRED: confirm funding sources, if any.
\textcolor{red}{[AUTHOR INPUT REQUIRED: confirm funding sources, if any.]}

\bmhead{Competing interests}
The authors declare no competing interests.

\bmhead{Ethics approval}
This study uses only anonymised, aggregated, publicly released third-party
mobility, meteorological and land-cover data. No data collection from human
participants was performed.

\bmhead{Availability of data and materials}
% ---------------------------------------------------------------------------
% Full catalogue of every real data source used in this project, as
% requested: name, provider, access point, and role. Kept here as the single
% authoritative list; update whenever a new source is added to the pipeline.
% ---------------------------------------------------------------------------
This study uses the following third-party data sources, all publicly
available:
\begin{itemize}
  \item \textbf{Road network}: OpenStreetMap contributors, regional extracts via
    Geofabrik GmbH (\url{https://download.geofabrik.de}), Open Database
    Licence.
  \item \textbf{Functional Urban Area boundaries}: JRC/OECD Global Human
    Settlement Layer, Functional Urban Areas \cite{ghsl_fua}.
  \item \textbf{Mobility demand}: MITMA (Ministerio de Transportes y
    Movilidad Sostenible, Spain) mobility study
    (\url{https://movilidad-opendata.mitma.es}).
  \item \textbf{Tree canopy / land cover}: ESA WorldCover 2021 v200
    \cite{esa_worldcover} (\url{https://esa-worldcover.org}), CC~BY~4.0.
  \item \textbf{Meteorological reanalysis and UTCI}: Copernicus Climate
    Change Service, ERA5-Land \cite{era5_land} and
    \texttt{derived-utci-historical-timeseries} \cite{cds_utci}, via the
    Copernicus Climate Data Store (\url{https://cds.climate.copernicus.eu}).
  \item \textbf{In-situ meteorological observations}: AEMET OpenData
    \cite{aemet_opendata} (\url{https://opendata.aemet.es}).
  \item \textbf{Socio-demographic vulnerability}: INE (Instituto Nacional de
    Estad\'istica) Atlas de Distribuci\'on de Renta de los Hogares
    \cite{ine_adrh} (\url{https://www.ine.es/dynt3/inebase/index.htm?padre=12385}),
    and census-section boundaries via INE's OGC API Features geometry
    service.
  \item \textbf{Land use / street trees / building height}: Copernicus Land
    Monitoring Service Urban Atlas \cite{urban_atlas}
    (\url{https://land.copernicus.eu/local/urban-atlas}), retrieved via the
    Copernicus Data Space Ecosystem.
  \item \textbf{Building/vegetation height}: CNIG/IGN PNOA second-coverage
    LiDAR-derived normalised digital surface models \cite{cnig_pnoa_lidar}
    (\url{https://centrodedescargas.cnig.es}).
\end{itemize}
Aggregate results and processed data tables supporting this study will be
made available upon publication.

\bmhead{Code availability}
The simulation wrapper, data-processing pipeline and analysis code
(\texttt{green\_mobility}) will be made available in a public repository upon
publication. The unmodified NOMAD simulation engine
\cite{nomad2026}, and the specific, narrowly-scoped, tested Python-binding
addition used to extract per-mode flow (Sec.~\ref{sec:flows}), are maintained
in a separate repository.

\bmhead{Author contributions}
% AUTHOR INPUT REQUIRED: confirm author contribution statement before submission.
\textcolor{red}{[AUTHOR INPUT REQUIRED: confirm author contribution
statement before submission.]}

\bibliography{references}

\end{document}
